% This is text associated with the open textbook "Open Electromagnetics"
% See immediately below for author, date
% This work is licensed under the Creative Commons Attribution-ShareAlike 4.0 International License. (CC BY-SA 4.0) 
% To view a copy of the license, visit http://creativecommons.org/licenses/by-sa/4.0/

%ORB TITLE Conductivity
%ORB AUTHOR 0        # S.W. Ellingson, Virginia Tech (ellingson@vt.edu)
%ORB DATE 20180806   # yyyymmdd
%ORB ABSTRACT Conductivity

\label{m0010_Conductivity} % <-- Must match name of this file and module directory name.  Don't mess with this!
\index{conductivity|(}

Conductivity is one of the three primary ``constitutive parameters'' that is commonly used to characterize the electromagnetic properties of materials (Section~\ref{m0007_Properties_of_Materials}).  The key idea is this:

%ORB HIGHLIGHT
\begin{mdframed}[backgroundcolor=yellow!20,nobreak=true] 
Conductivity 
\index{conductivity!definition}
is a property of materials that determines conduction current density in response to an applied electric field.
\end{mdframed}
%
Recall that conduction current is the flow of charge in response to an electric field (Section~\ref{m0110_Convection_and_Conduction}).
Although the associated force is straightforward to calculate (e.g., Section~\ref{m0102_Coulombs_Law}), the result is merely the force applied, not the speed at which charge moves in response.  The latter is determined by the \emph{mobility of charge}, 
\index{charge mobility}
which is in turn determined by the atomic and molecular structure of the material.  
Conductivity relates current density to the applied field directly, without requiring one to grapple separately with the issues of applied force and charge mobility.

In the absence of material -- that is, in a true, perfect vacuum -- conductivity is zero because there is no charge available to form current, and therefore the current is zero no matter what electric field is applied.   
At the other extreme, a \emph{good conductor} 
\index{conductor!good}
is a material that contains a supply of charge that is able to move freely within the material.  
When an electric field is applied to a good conductor, charge-bearing material constituents move in the direction determined by the electric field, creating current flow in that direction.  
This relationship is summarized by \emph{Ohm's Law for Electromagnetics}:
\index{Ohm's law}
%
\begin{equation}
\boxed{
{\bf J} = \sigma {\bf E}
}
\label{m0010_eMatOhmsLaw}
\end{equation}
%
where 
${\bf E}$ is electric field intensity (V/m);
\index{electric field intensity!related to current}
${\bf J}$ is the volume current density, a vector describing the current flow, having units of A/m$^2$ (see Section~\ref{m0101_Current_Distributions}); and
$\sigma$ is conductivity.
Since ${\bf E}$ has units of V/m, we see $\sigma$ has units of $\Omega^{-1}\mbox{m}^{-1}$, which is more commonly expressed as S/m, where 1~S (``siemens'') is defined as $1~\Omega^{-1}$.  
Section~\ref{m0137_Table_of_Conductivities} provides values of conductivity for a representative set of materials. 
%
%ORB HIGHLIGHT
\begin{mdframed}[backgroundcolor=yellow!20,nobreak=true] 
Conductivity $\sigma$ is expressed in units of S/m, where 1~S $=$ 1~$\Omega^{-1}$.  
\end{mdframed}
%
It is important to note that the current being addressed here is \emph{conduction current}, and not convection current, displacement current, or some other form of current -- see Section~\ref{m0101_Current_Distributions} for elaboration.  Summarizing:
%
%ORB HIGHLIGHT
\begin{mdframed}[backgroundcolor=yellow!20,nobreak=true] 
Ohm's Law for Electromagnetics (Equation~\ref{m0010_eMatOhmsLaw}) states that volume density of conduction current (A/m$^2$) equals conductivity (S/m) times electric field intensity (V/m). 
\end{mdframed}

The reader is likely aware that there is also an ``Ohm's Law'' in electric circuit theory, which states that current $I$ (units of A) is voltage $V$ (units of V) divided by resistance $R$ (units of $\Omega$); i.e., $I=V/R$.  This is in fact a special case of Equation~\ref{m0010_eMatOhmsLaw}; see Section~\ref{m0071_Resistance_DC} for more about this. 

As mentioned above, $\sigma$ depends on both the \emph{availability} and \emph{mobility} of charge within the material.  
At the two extremes, we have \emph{perfect insulators}, 
\index{insulator!perfect}
for which $\sigma=0$, and \emph{perfect conductors}, 
\index{conductor!perfect}
for which $\sigma \to \infty$.  
Some materials approach these extremes, whereas others fall midway between these conditions.  
Here are a few classes of materials that are frequently encountered:
%
\begin{itemize}

\item A perfect vacuum -- ``free space'' -- contains no charge and therefore is a perfect insulator with $\sigma=0$. 

\item \emph{Good insulators} typically have conductivities $\ll 10^{-10}$~S/m, which is sufficiently low that the resulting currents can usually be ignored.  The most important example is air, which has a conductivity only slightly greater than that of free space.  An important class of good insulators is {\it lossless dielectrics},\footnote{See Section~\ref{m0107_Dielectric_Media} 
\index{dielectric (material)!lossless}
for a discussion of dielectric materials.} which are well-characterized in terms of permittivity
\index{permittivity}
($\epsilon$) alone, and for which $\mu=\mu_0$ and $\sigma=0$ may usually be assumed. 

\item \emph{Poor insulators} 
\index{insulator!poor}
have conductivities that are low, but nevertheless sufficiently high that the resulting currents cannot be ignored.  For example, the dielectric material that is used to separate the conductors in a transmission line must be considered a poor insulator as opposed to a good (effectively lossless) insulator in order to characterize loss per length along the transmission line, which can be significant.\footnote{Review Sections~\ref{m0029_Lumped-Element_Model}, \ref{m0083_Lossless_Transmission_Lines}, and associated sections for a refresher on this issue.}
These \emph{lossy dielectrics} 
\index{dielectric (material)!lossy}
are well-characterized in terms of $\epsilon$ and $\sigma$, and typically $\mu=\mu_0$ can be assumed. 

\item \emph{Semiconductors} 
\index{semiconductor (material)}
such as those materials used in integrated circuits have intermediate conductivities, typically in the range $10^{-4}$ to $10^{+1}$~S/m.

\item \emph{Good conductors} 
\index{conductor!good}
are materials with very high conductivities, typically greater than $10^5$~S/m.  
An important category of good conductors includes metals, with certain metals including alloys of aluminum, copper, and gold reaching conductivities on the order of $10^8$~S/m. 
In such materials, minuscule electric fields give rise to large currents.   
There is no significant storage of energy in such materials, and so the concept of permittivity is not relevant for good conductors. 

\end{itemize}

The reader should take care to note that terms such as {\it good conductor} and {\it poor insulator} are qualitative and subject to context.  What may be considered a ``good insulator'' in one application may be considered to be a ``poor insulator'' in another.  

One relevant category of material was not included in the above list -- namely, \emph{perfect conductors}. 
A perfect conductor is a material in which $\sigma\to\infty$.  
It is tempting to interpret this as meaning that any electric field gives rise to infinite current density; however, this is not plausible even in the ideal limit.
Instead, this condition is interpreted as meaning that ${\bf E}={\bf J}/\sigma\to 0$ throughout the material; i.e., ${\bf E}$ is zero independently of any current flow in the material.
An important consequence is that the potential field $V$ is equal to a constant value throughout the material.  
(Recall ${\bf E}=-\nabla V$ (Section~\ref{m0063_Electric_Field_as_the_Gradient_of_Potential}), so constant $V$ means ${\bf E}=0$.)
We refer to the volume occupied by such a material as an \emph{equipotential volume}.
This concept is useful as an approximation of the behavior of good conductors; for example, metals are often modeled as perfectly-conducting equipotential volumes in order to simplify analysis.
%
%ORB HIGHLIGHT
\begin{mdframed}[backgroundcolor=yellow!20,nobreak=true] 
A perfect conductor 
\index{conductor!perfect (definition)}
is a material for which $\sigma\to\infty$, ${\bf E}\to 0$, and subsequently $V$ (the electric potential) is constant.  
\end{mdframed}

One final note: It is important to remain aware of the assumptions we have made about materials in this book, which are summarized in Section~\ref{m0007_Properties_of_Materials}. In particular, we continue to assume that materials are linear unless otherwise indicated.  
For example, whereas air is normally considered to be a good insulator and therefore a poor conductor, anyone who has ever witnessed lightning 
\index{lightning}
has seen a demonstration that under the right conditions -- i.e., sufficiently large potential difference between earth and sky -- current will readily flow through air.  This particular situation is known as \emph{dielectric breakdown} 
\index{dielectric breakdown}
(see Section~\ref{m0109_Dielectric_Breakdown}).  The non-linearity of materials can become evident before reaching the point of dielectric breakdown, so one must be careful to consider this possibility when dealing with strong electric fields.

{\bf Additional Reading:}
\begin{itemize}
\item ``\href{https://en.wikipedia.org/wiki/Electrical_resistivity_and_conductivity}{Electrical Resistivity and Conductivity}'' on Wikipedia.
\item ``\href{https://en.wikipedia.org/wiki/Ohm's_law}{Ohm's Law}'' on Wikipedia.
\end{itemize}

\index{conductivity|)}
